% Generated draft: requires mathematical and source review.
\documentclass[11pt]{article}
\usepackage[margin=25mm]{geometry}
\usepackage{fontspec}
\setmainfont{Latin Modern Roman}
\usepackage{amsmath,amssymb,mathtools,unicode-math}
\setmathfont{Latin Modern Math}
\usepackage{enumitem,needspace,xurl,hyperref}
\setlength{\parindent}{0pt}
\setlength{\parskip}{5pt}
\setlength{\emergencystretch}{3em}
\newcommand{\N}{\mathbb{N}}
\newcommand{\Z}{\mathbb{Z}}
\newcommand{\Q}{\mathbb{Q}}
\newcommand{\R}{\mathbb{R}}
\newcommand{\C}{\mathbb{C}}
\newcommand{\F}{\mathbb{F}}
\newcommand{\A}{\mathbb{A}}
\newcommand{\PP}{\mathbb{P}}
\newcommand{\E}{\mathbb{E}}
\newcommand{\eps}{\varepsilon}
\newcommand{\dd}{\,\mathrm{d}}
\begin{document}
\section*{UnsolvedMath ID: 30002491}
\textbf{Source problem code:} catalogue entry
\subsection{Definitions and mathematical statement}
A local system has a monodromy representation once an appropriate fundamental group, coefficient field, and fiber have been specified. Schematically, write this representation as
\[
\rho_{\mathcal L}\colon \Pi\longrightarrow \operatorname{GL}_K(V),
\]
where \(K\) is a field and \(V\) is a finite-dimensional \(K\)-vector space. Its algebraic monodromy group is the Zariski closure
\[
G_{\mathcal L}=\overline{\rho_{\mathcal L}(\Pi)}^{\mathrm{Zar}}\subseteq\operatorname{GL}_K(V).
\]
These symbols describe the invariant under discussion; the supplied record does not fix \(K\), the base space, the category of local systems, or the fundamental group \(\Pi\). In particular, it does not specify arithmetic versus geometric monodromy.

Exceptional simple algebraic group types are \(G_2,F_4,E_6,E_7,E_8\). The record does not specify which types or forms are intended, or whether exceptional monodromy refers to the full algebraic monodromy group, its identity component, or a related derived group.

\textbf{Source question.} Can the automorphic and constructive approaches to producing local systems with prescribed exceptional algebraic monodromy be brought together?

A schematic expression of the prescribed invariant is \(G_{\mathcal L}\cong G\), for an intended exceptional algebraic group \(G\); this expression does not supply the missing hypotheses or quantifiers. The source gives neither a definition of the two approaches nor a mathematical criterion for their unification. Consequently, the question is retained as a research program: the supplied evidence does not determine a precise universally or existentially quantified conjecture. In particular, it does not justify requiring constructions for every exceptional group or asserting that two specified local systems are isomorphic.
\subsection{Short English statement}
The question asks whether automorphic methods and constructive methods for producing local systems with prescribed exceptional algebraic monodromy can be unified. The setting, scope, and meaning of unification remain unspecified, so this is a statement-only research question requiring mathematical review.
\subsection{Sources}
\begin{itemize}
\item Supplied local UnsolvedMath registry record, \texttt{data/unsolvedmath/problems.json}, permanent ID \(30002491\), exact source problem code \texttt{catalogue entry}; title: \emph{Unifying Constructions of Exceptional Monodromy Local Systems}. The record gives proposed year 2014 and names no individual proposer.
\item Supplied report citation: \emph{Langlands Correspondence and Constructive Galois Theory} (2014), Oberwolfach Reports, pp.\ 287--334, DOI 10.4171/owr/2014/05. \url{https://doi.org/10.4171/owr/2014/05}. The supplied background describes Oberwolfach Reports as documenting mathematical workshops and their open-problem discussions; no precise page or individual contribution for this question is supplied.
\item Supplied literature source: \emph{Kloosterman sheaves for reductive groups}. \url{https://arxiv.org/abs/1005.2765}.
\item Literature assessment dated \texttt{2026-08-21}: the supplied curation labels the question open and states that later work constructs important local systems and exceptional monodromy examples, but that no general synthesis or explicit trace calculation answers this broad question. The listed sources checked are the OWR report and \emph{Kloosterman sheaves for reductive groups}, linked above. This is supplied dated curation, not an independently verified literature assessment; its time-sensitive status and cited claims require independent verification.
\item The supplied registry labels the statement \texttt{corrected\_verified} and says its mathematical content was unchanged while notation, typography, and wording were normalized into a self-contained publication form. This records the supplied editorial assessment, not independent mathematical verification of this draft.
\end{itemize}
\end{document}
